\section{The literal matched coefficient and its distinguished zero mode}
\label{sec:literal-zero-mode}

We restate the minimum part of the TPC-32 packet needed for the
normalization analysis.  This prevents an abstract coefficient
sequence from being silently substituted for the physical object.

\subsection{Rows, targets, and complete matched atoms}

Let the physical rows be
\[
 \alpha=(\ell_\alpha,d_\alpha),\qquad
 m_\alpha=\ell_\alpha d_\alpha\asymp Q,
\]
and similarly for \(\gamma\).  For an allowed orbit index \(j\), set
\begin{equation}
 N_\alpha(j)=m_\alpha j+h_0,
 \qquad
 N_\gamma(j)=m_\gamma j+h_0.
 \label{eq:literal-targets}
\end{equation}
The primitive support conditions inherited from the packet ensure
that every common target divisor can be cancelled against \(j\), so
that it divides \(m_\alpha-m_\gamma\)
\citep{WangTPC30,WangTPC32}.

For a cutoff \(Y\), write the cumulative divisor prefix as
\begin{equation}
 \mathsf A_{m,Y}(j)
 =\sum_{\substack{w\le Y\\w\mid mj+h_0}}b_R(w),
 \label{eq:literal-divisor-prefix}
\end{equation}
where \(b_R\) is the divisor coefficient fixed by the physical
construction.  The matched cutoff kernel is
\begin{equation}
 \cK^{\rm sh}_{\alpha,\gamma}(j)
 =
 \mathsf A_{m_\alpha,U_0}(j)
 \mathsf A_{m_\gamma,U_0}(j)
 -
 \mathsf A_{m_\alpha,T}(j)
 \mathsf A_{m_\gamma,T}(j).
 \label{eq:literal-matched-kernel}
\end{equation}
Put
\[
 \mathsf C_m
 =\mathsf A_{m,U_0}-\mathsf A_{m,T},
 \qquad
 \mathsf H_m
 =\frac{\mathsf A_{m,U_0}+\mathsf A_{m,T}}2
 =\mathsf A_{m,T}+\frac12\mathsf C_m.
\]
The three physical channels and the exact two-leg polarization are
\begin{equation}
 \cK^{\rm sh}_{\alpha,\gamma}
 =
 \mathsf A_{m_\alpha,T}\mathsf C_{m_\gamma}
 +\mathsf C_{m_\alpha}\mathsf A_{m_\gamma,T}
 +\mathsf C_{m_\alpha}\mathsf C_{m_\gamma}
 \label{eq:literal-three-channel-expansion}
\end{equation}
and
\begin{equation}
 \boxed{
 \cK^{\rm sh}_{\alpha,\gamma}
 =\mathsf C_{m_\alpha}\mathsf H_{m_\gamma}
  +\mathsf H_{m_\alpha}\mathsf C_{m_\gamma}.}
 \label{eq:literal-two-leg-polarization}
\end{equation}
Both terms in \eqref{eq:literal-two-leg-polarization} are retained;
none of the arguments below deletes a polarization or separates the
three channels by a triangle inequality.

Let \(\cI_X\) denote the complete finite set of allowed triples
\(\omega=(\alpha,\gamma,j)\).  Define
\begin{equation}
 w_X(\omega)
 =
 \gamma_\alpha^{(1)}\gamma_\gamma^{(2)}
 \mathfrak A_{\alpha,\gamma}(j)
 \cK^{\rm sh}_{\alpha,\gamma}(j).
 \label{eq:literal-complete-atom}
\end{equation}
Here the two \(\gamma\)'s are the literal opened-row coefficients and
\(\mathfrak A_{\alpha,\gamma}(j)\) is the actual joint multiplier.  It
contains the off-diagonal row mask, the fixed residue conditions, and
the physical smooth weights.  Thus \(w_X\) is a signed or complex
amplitude, not its absolute majorant.

\begin{remark}[Calibration convention]
\label{rem:literal-calibration}
The calibrated prefix used in TPC-32 moves an explicit mean term into
two drift legs.  Those legs are part of the literal matched-shell
ledger.  The operation does not assert that the determinant
coefficient introduced below has zero total sum.
\end{remark}

\subsection{Content and normalized determinant}

For \(\omega=(\alpha,\gamma,j)\), put
\begin{equation}
 c(\omega)
 =\gcd(N_\alpha(j),N_\gamma(j)),
 \qquad
 \Delta^\#(\omega)
 =\frac{m_\alpha-m_\gamma}{c(\omega)}.
 \label{eq:literal-content-determinant}
\end{equation}
The primitive support makes \(\Delta^\#(\omega)\) an integer whenever
the off-diagonal atom is active.  At the physical cutoff
\begin{equation}
 \boxed{C=\lfloor J\rfloor,}
 \label{eq:literal-C-equals-J}
\end{equation}
define
\begin{equation}
 \boxed{
 A_C(n)
 =
 \sum_{\omega\in\cI_X}
 w_X(\omega)
 \one_{c(\omega)\le C}
 \one_{\Delta^\#(\omega)=n},
 \qquad n\in\Z.}
 \label{eq:literal-AC}
\end{equation}
This is aggregation before squaring.  The two physical quantities of
interest are
\begin{equation}
 \boxed{
 Z_X=\sum_{n\in\Z}A_C(n),
 \qquad
 D_X=\sum_{n\in\Z}|A_C(n)|^2.}
 \label{eq:literal-Z-D}
\end{equation}
The first is the distinguished zero mode; the second is the literal
post-bin energy.

\begin{proposition}[Two meanings of zero]
\label{prop:two-zeros}
The off-diagonal row mask implies
\begin{equation}
 A_C(0)=0.
 \label{eq:literal-bin-zero}
\end{equation}
For every auxiliary modulus \(q\), however,
\begin{equation}
 \widehat A_{C,q}(0)
 :=\sum_{n\in\Z}A_C(n)e_q(0)
 =Z_X,
 \label{eq:literal-fourier-zero}
\end{equation}
and \eqref{eq:literal-bin-zero} gives no cancellation in
\eqref{eq:literal-fourier-zero}.
\end{proposition}

\begin{proof}
If \(\Delta^\#(\omega)=0\), then
\(m_\alpha=m_\gamma\).  Such an atom is deleted by the actual
off-diagonal mask, proving \eqref{eq:literal-bin-zero}.  At Fourier
frequency \(r=0\), the character
\(e_q(-r n)\) equals one for every determinant bin \(n\), which gives
\eqref{eq:literal-fourier-zero}.  Deleting the single bin \(n=0\)
does not force the sum over the remaining bins to vanish.
\end{proof}

\begin{remark}[Frequency typing]
The \(r=0\) in \eqref{eq:literal-fourier-zero} is the finite Fourier
frequency of the normalized determinant.  It is not an orbit-variable
Poisson zero mode and not a zero mode of a centered divisor kernel.
\end{remark}

\subsection{The full-shell splice}

Let
\begin{align}
 \cS_{\rm full}^{\rm sh}
 &=\sum_{\omega\in\cI_X}w_X(\omega),\\
 \cS_{>C}^{\rm sh}
 &=\sum_{\omega\in\cI_X}
   w_X(\omega)\one_{c(\omega)>C}.
 \label{eq:literal-large-content-shell}
\end{align}

\begin{proposition}[Exact zero-mode/full-shell identity]
\label{prop:zero-full-shell}
At the literal cutoff \eqref{eq:literal-C-equals-J},
\begin{equation}
 \boxed{
 Z_X
 =\cS_{\rm full}^{\rm sh}-\cS_{>C}^{\rm sh}.}
 \label{eq:literal-zero-full-shell}
\end{equation}
In the inherited physical range,
\begin{equation}
 |\cS_{>C}^{\rm sh}|
 \ll_{\varepsilon,h_0,\cD}
 X^\varepsilon
 \left(Q^2+\frac{XQ}{C}\right)
 =X^{o(1)}Q^2.
 \label{eq:literal-large-content-bound}
\end{equation}
Consequently, up to an \(X^{o(1)}Q^2\) error, estimating \(Z_X\) is
the same problem as estimating the complete matched shell.
\end{proposition}

\begin{proof}
Exact contents partition \(\cI_X\), so summing
\eqref{eq:literal-AC} over \(n\) gives the small-content shell.  Its
complement is \eqref{eq:literal-large-content-shell}, proving
\eqref{eq:literal-zero-full-shell}.  The first inequality in
\eqref{eq:literal-large-content-bound} is the inherited
large-content estimate \citep{WangTPC30,WangTPC32}.  The physical
scale relation \(JQ=X^{1+o(1)}\), together with
\(C=\lfloor J\rfloor\), gives
\(XQ/C=X^{o(1)}Q^2\).
\end{proof}

\begin{corollary}[The coherent completion required from generic
frequencies]
\label{cor:generic-coherent-completion}
Under the strong physical no-wrap padding
\eqref{eq:constant-no-wrap-q}, the distinguished mode satisfies
\[
 Z_X=-\sum_{r\ne0}\widehat A_{C,q}(r).
\]
Consequently, a generic-frequency or mean-square theorem controls
\(Z_X\) only after it also controls the coherent sum of the retained
nonzero modes and every exceptional nonzero mode.  A direct estimate
at \(r=0\) is sufficient, but it is not logically necessary.
\end{corollary}

\begin{proof}
This is \eqref{eq:constant-coherent-reassembly}.  Bounds for most
individual modes do not estimate the displayed sum when the
uncontrolled nonzero modes or the phase-coherent reassembly can carry
the full mass.  Conversely, any estimate for the complete coherent
sum is already an estimate for \(Z_X\), regardless of whether a
separate pointwise theorem was stated at \(r=0\).
\end{proof}
